\scanpage{190}\mspage{46}
% archive sheet number: 190

$\mathrm{F}\ 0)$ $\pi_0(\unsure{f|I|u})$ est bijectif

$\mathrm{F}\ 1)$ Pour tt $i\in\mathrm{Ob}\,I$, existe $i'\in I'$ et une flèche $i\longrightarrow u(i')$ de $I$.

% slanted
\margin{\uline{NB} F(1) implique déjà surjectivité \unsure{du|sur} $\pi_0$, donc sa bijectivité si $I'$ est $0$-connexe -- dans ce cas F0 \deleted{\unsure{est}} résulte de F1)}

% slanted
\margin{Si $I'$ est loc. \unsure{fblt} filtr. (= préfiltrante), alors pour la cofinalité de $u$, il \uline{suffit} qu'on ait F(0), F(1), F(2), et alors les $\,{}_{i\backslash}I'$ \unsure{sont} fblt filtrantes, et $I$ est préfiltrant tout comme $I'$ !}

b) Si $I'$ est fblt filtrante, pour que $u$ soit cofinal, il \uline{suffit} qu'elle satisfasse F1, ainsi que la condition :

$\mathrm{F}\ 2)$ Pour toute double flèche

\diagram{\begin{tikzcd}[ampersand replacement=\&]
i \arrow[r,shift left=.6ex,"\beta"] \arrow[r,shift right=.6ex,"\gamma"'] \& u(i')
\end{tikzcd}}

\deleted{existe} (avec $i\in\mathrm{Ob}\,I$, $i'\in\mathrm{Ob}\,I$) existe une flèche $\deleted{\unsure{b'}}i'\xrightarrow{\alpha'}j'$ de $I'$, de telle façon que $u(\alpha')$ égalise $\beta,\gamma$, i.e. $u(\alpha')\beta=u(\alpha')\gamma$.

Et ces conditions F1, F2 impliquent que les $\,{}_{i\backslash}I'$ sont fblt filtrantes, et $I$ aussi.

% slanted
\margin{Si $I'$ est loc filtr.-te il est cofinal ssi \deleted{\unsure{elle}} satisfait \deleted{\unsure{alors}} F0, F1, F2, et \deleted{\unsure{alors}} les $\,{}_{i\backslash}I'$ sont filtr.-tes, et alors $I$ est loc. filtrante.}

c) Si $I'$ est filtrante, alors $u$ est cofinal \uuline{ssi} il satisfait les deux conditions F1, F2, ou \deleted{\unsure{est}} encore ssi les $\,{}_{i\backslash}I'$ sont \deleted{\unsure{fblt}} \uline{filtrantes} (et non seulement $0$-connexes). Et cela implique que $\unsure{I|I'}$ est \inserted{également} filtrante.

d) Si $I$ est \inserted{\unsure{pré}}filtrante \inserted{(= loc. \unsure{fblt.|fbt.} filtrante, i.e. PS.1)} \inserted{et $u$} \uline{pl. fidèle}, pour que $u$ soit cofinal, il \uline{f et s} que $u$ satisfasse la condition F1 plus haut. Et alors $I'$ est également \deleted{filtrante,} \deleted{loc. filtrante} préfiltrante, et les $\,{}_{i\backslash}I'$ sont fblt. filtrantes. Si de plus $I$ est \deleted{connexe, i.e. fblt filtrante} \deleted{\unsure{(resp.} si $I$ est} filtrante\deleted{\unsure{s}}, \inserted{les $\,{}_{i\backslash}I'$ et} $I'$ \unsure{l'est|l'ét} également.

(Comparer SGA 4 I 8.1.3.)
